%! @@TITLE@@ — Unit @@UNITINT@@, Lesson @@LESSONID@@ — Lesson Plan
% GRADUAL-RELEASE plan: warm-up → guided notes & practice → group activity → debrief → close.
% Teacher-facing; never handed to students.
\documentclass[10pt]{article}
\usepackage{@@PREFIX@@-article}
\usepackage{@@PREFIX@@-boxes}
\usepackage{graphicx}   % -article does NOT load graphicx; needed for thumbnails/screenshots
\graphicspath{{images/}}
\newcommand{\TallMath}[1]{$\displaystyle #1\rule[-1.4em]{0pt}{3.2em}$}
@@COURSEMACROS@@\newcommand{\UnitNumberName}{Unit @@UNITINT@@: @@UNITTITLE@@ \quad}
\newcommand{\LessonNumberName}{Lesson @@LESSONID@@: @@TITLE@@}

\begin{document}
\begin{center}
    {\Large\bfseries \CourseName \\
    {\small\bfseries \UnitNumberName \LessonNumberName}}
\end{center}

\begin{tcolorbox}[colback=forestbg, colframe=forest, boxrule=0.9pt, arc=2mm,
    left=2mm, right=2mm, top=1.5mm, bottom=1.5mm]
    \textbf{Primary Objective:} TODO --- what students can do/interpret/justify by the end.
    \\[2pt]
    \textbf{Standards (2023 VA SOL):} TODO --- the codes the user supplied, with clause notes.
    \\[2pt]
    \textbf{Lesson model:} \emph{Gradual release} --- I do, we do, you do together, you do alone. A
    warm-up activates prior knowledge; \textbf{Guided Notes} build and name the vocabulary with the
    class; a \textbf{Group Activity} applies it to new contexts and a model; a \textbf{debrief}
    consolidates and surfaces errors; \textbf{homework} is the independent practice.
\end{tcolorbox}

\renewcommand{\arraystretch}{1.4}
\begin{skillbox}[Priority Ideas \& Skills]{goldbox}
% Fill in the \item text ONLY. Two tabularx cells, one itemize each — do not re-nest or
% re-balance the environments (that is the classic source of a stray \end{itemize}).
\begin{tabularx}{\linewidth}{@{}X|X@{}}
\textbf{Priority Ideas \& Skills} & \textbf{Key Understandings (the ``why'')} \\[2pt]
\hline
\vspace{-6pt}
\begin{itemize}[leftmargin=1.1em, nosep, after=\vspace{2pt}]
  \item TODO: a skill, stated as something the student does.
\end{itemize}
&
\vspace{-6pt}
\begin{itemize}[leftmargin=1.1em, nosep, after=\vspace{2pt}]
  \item TODO: the why --- include the lesson's \textbf{target misconception} stated explicitly.
\end{itemize}
\end{tabularx}
\end{skillbox}

\begin{skillbox}[{Vocabulary, Concepts, \& Theorems}]{sky}
    % TODO: key terms + definitions (these are the terms the Guided Notes vocabbox builds).
\end{skillbox}

\begin{fixedskillbox}[Lesson at a Glance --- \MeetingLength]{forestbg}
\renewcommand{\arraystretch}{1.25}
\begin{tabularx}{\linewidth}{@{}l c X X@{}}
\textbf{Phase} & \textbf{Min} & \textbf{Students} & \textbf{Teacher} \\
\hline
Warm-Up & 5 & TODO: three prior skills, alone then check with a neighbour & Circulate; debrief the one item that seeds the notes \\
Guided Notes & 15 & Fill the vocabulary and the four sections with the class & \emph{I do / we do}: build each name on the hook's numbers; work the \texttt{work} blocks at the board \\
Group Activity & 25 & \emph{TODO Title} parts 1--4 in groups of 3--4 & Circulate; questions, cues, prompts; choose work to display \\
Debrief & 10 & Share findings; correct their own pages & Put the crux item and the comparison on the board \\
Close \& Homework & 5 & Record the assignment; preview the next lesson & Assign the homework page; name what changed today \\
\end{tabularx}
\end{fixedskillbox}

\begin{fixedskillbox}[Warm-Up --- Activate Prior Knowledge \& Spiral Review]{forestbg}
    % TODO: the ~3 warm-up items, which prior skill each rehearses, and what to debrief aloud.
    % Carry the debriefed observation straight into notes section 1 — that is the handoff.
@@SPIRALWARMUP@@
\end{fixedskillbox}

\begin{skillbox}[Hook]{forestbg}
    TODO: the 60-second context, the 3--4 questions to ask with their answers in parentheses, and
    the idea to land in one sentence. The hook's numbers should be the ones every worked example in
    the Guided Notes reuses --- that is what makes the notes read as one lesson.
\end{skillbox}

\begin{skillbox}[Guided Notes --- the Lesson (15 min)]{forestbg}
\begin{multicols}{2}
\textbf{1. TODO.} TODO --- what to build, and the first \texttt{work} block.

\textbf{2. TODO.} TODO.

\columnbreak
\textbf{3. TODO.} TODO --- this is normally the section carrying the target misconception; say so.

\textbf{4. TODO.} TODO --- the special case and the boundary.

\textbf{Practice (the ``we do'').} TODO: the guided-practice example. Do item~1 together, then
release the rest and circulate.
\end{multicols}
\end{skillbox}

\begin{skillbox}[Group Work --- \emph{TODO Title} (25 min, groups of 3--4)]{redbox}
One common task, \textbf{no tiers}. \textbf{Launch (2 min):} TODO the launch script. One packet
each, one shared conversation, and the rule: \emph{for every answer, be ready to show me where you
see it.}
\begin{multicols}{2}
\textbf{What students do.} TODO: the arc of parts 1--4 --- what part~1 establishes, how part~2
varies it and where the \textbf{crux question} sits, what part~3 compares, and what the part~4
model asks.

\columnbreak
\textbf{What the teacher does.} Circulate. Listen before speaking. Give a \emph{question, cue, or
prompt}, not an answer:
\begin{itemize}[leftmargin=1.1em, nosep]
  \item TODO: item ref: ``TODO prompt.''
\end{itemize}
Pick the work you will display in the debrief.
\end{multicols}
\end{skillbox}

\boxguard
\begin{skillbox}[Debrief (10 min --- what goes on the board, in a second color)]{forestbg}
Do not re-teach the activity; land four things on the displayed student work.
\begin{enumerate}[leftmargin=1.3em, nosep, itemsep=2pt]
  \item \textbf{TODO: the model answer.} Label a group's part~1 completely; everyone corrects
        their own page as you go.
  \item \textbf{TODO: the crux.} Use a group's wording verbatim. This is the misconception the
        unit test will probe --- say it twice.
  \item \textbf{TODO: the comparison} from part~3, and the general conclusion it licenses.
  \item \textbf{TODO: the model's concept check} from part~4, and what a wrong answer reveals.
\end{enumerate}
TODO: which item to cut if time is short, and which two to protect.
\end{skillbox}

\boxguard
\begin{skillbox}[Active Monitoring --- Watch For]{redbox}
\begin{itemize}[leftmargin=1.2em, nosep]
  \item TODO: misconception, keyed to an activity or homework item number.
\end{itemize}
Cold-call prompts: TODO.
\end{skillbox}

\boxguard
\begin{skillbox}[Reinforcement \& Extension]{goldbox}
\textbf{Homework} (in the packet, collected next class): TODO --- itemize the ~6 problems and the
extension. \textbf{DeltaMath override:} TODO --- say whether this content is well covered in
DeltaMath and, if so, what set to swap in instead of the packet page. \textbf{Preview:}
TODO next lesson.
\end{skillbox}

% ── Teacher notes — one per component, in packet order ──
% Four notes: Warm-Up, Guided Notes, Group Activity, Homework. This is the ONLY place teacher
% prose goes — never append one to a _key, which would make the key longer than its blank.
\begin{teachernote}[Warm-Up]
TODO: what each item seeds, what to debrief, what to cut if it runs long.
\end{teachernote}

\begin{teachernote}[Guided Notes]
TODO: the per-section pacing split that fills 15 minutes; which section is the one that matters;
where the \texttt{work} blocks are; what to cut if you fall behind.
\end{teachernote}

\begin{teachernote}[Group Activity]
TODO: the per-part time budget that fills 25 minutes; the common slips; the early-finisher move.
\end{teachernote}

\begin{teachernote}[Homework]
TODO: which item separates students and why; the formative check (the multiple-choice item) and
the categories to sort responses into; how many \texttt{work} blocks the page carries.
\end{teachernote}

\end{document}
